\section{Related Work}












\subsection{Training-Dependent Conversion}
 
Training-dependent methods retrain the source ANN to align activations with spiking neuron dynamics before conversion. Under rate coding, representative works optimize membrane potential initialization~\cite{bu2022optimized}, replace ReLU with quantization-aware activations~\cite{bu2023optimal}, or apply dual-phase pipelines for weight and threshold co-calibration~\cite{ding2021optimal, wang2022towards}. For Transformers, SpikeZIP-TF~\cite{you2024spikezip} fine-tunes pretrained weights with spike-compatible surrogates, and SpikeLLM~\cite{xing2024spikellm} extends spiking mechanisms to billion-parameter LLMs via saliency-aware encoding. While effective, these methods require training pipeline access, limiting applicability when the source ANN is fixed or proprietary.

\subsection{Training-Free Conversion}
 
In contrast, Training-free approaches convert pretrained ANNs directly without fine-tuning. The dominant paradigm is rate coding~\cite{rueckauer2017conversion, diehl2015fast}, which requires at least $T = N$ timesteps to distinguish $N$ activation levels. To mitigate this, prior works calibrate residual and initial membrane potentials~\cite{hao2023reducing, hao2023bridging}, extend neuron outputs to ternary values~\cite{wang2022signed}, or refine threshold calibration strategies~\cite{wang2025adaptive, bu2025inference}. Beyond rate coding, phase coding~\cite{kim2018deep, wang2022efficient} supports $2^T$ representable values but requires all timesteps to elapse before decoding, while differential coding (DCGS)~\cite{huang2025differential} eliminates temporal decay at the cost of hardware-unfriendly division operations. For Transformer conversion, ECMT~\cite{huang2024towards} employs multi-threshold neurons but retains floating-point multiplication; STA~\cite{jiang2024spatio} approximates arbitrary nonlinearities but requires hundreds of timesteps; and SpikedAttention~\cite{hwang2024spikedattention} achieves fully spike-driven conversion via trace-driven matrix multiplication. 
Despite these advances, all existing methods rely on iterative accumulation and gradual dissipation of residual membrane potential, encoding redundant information across timesteps — a fundamental bottleneck that neither algorithmic nor hardware optimizations have yet resolved.

In contrast, PRIME belongs to the training-free category and requires no fine-tuning or retraining: the conversion is performed in a single forward pass on the pretrained weights, directly parsing the binary representation of floating-point data to generate spikes.